\tzpanchor{0705}
\tzpentryheader{0705}{BERRY CONNECTION GAUGE INVARIANCE}

\begingroup\small
\begingroup\scriptsize
\setlength{\tabcolsep}{4pt}
\renewcommand{\arraystretch}{0.92}
\noindent\begin{tabularx}{\linewidth}{@{}p{0.24\linewidth}p{0.36\linewidth}X@{}}
\textbf{UUID} & \multicolumn{2}{l@{}}{\tzpcode{4508e95b-314b-5fdc-b54f-f21335b340b8}}\\
\textbf{PiID} & \tzpcode{9561121290219} & \textbf{TZPi}\quad \tzpcode{1}\quad \textbf{PiPE}\quad \tzpcode{712}\\
\textbf{RTEID} & \textbf{Poloidal $\theta$}\quad \tzpcode{2706.6058 rad} & \textbf{Toroidal $\phi$}\quad \tzpcode{04.4597 rad}\\
\textbf{TZPID} & \tzpcode{ID0705} & \textbf{Node Dependencies}\quad \tzpidref{0704}\\
\textbf{Registry} & \tzpcode{registry\_id\_0705} & \textbf{Scale}\quad transfer\\
\textbf{LOC Classification} & QA641 manifolds and topology & QC174.17 mathematical physics\\
\textbf{arXiv Classification} & Primary: math-ph & Cross-list: gr-qc, math.DG\\
\textbf{Primary Support} & Berry Connection & Geometric Phase\\
\textbf{Closure Support} & Gauge Shift & \\
\end{tabularx}\par
\endgroup
\endgroup

\tzpsection{Canonical Statement}
The entry records berry connection gauge invariance through a normalized Octonary TRAWIN operator chain.

\tzpsection{Canonical Equation}
\[
\mathcal{O}_{\mathrm{id0705}}:\mathcal{X}_{\mathrm{id0705}}\longrightarrow\mathcal{Y}_{\mathrm{id0705}}
\]
\[
\mathcal{N}_{\mathrm{TRAWIN}}(\mathcal{O}_{\mathrm{id0705}})=\mathcal{O}_{\mathrm{id0705}}
\]
\[
\mathcal{O}_{\mathrm{id0705}}(x)\in\mathcal{Y}_{\mathrm{adm}}
\]

\begin{minipage}[t]{0.49\linewidth}
\tzpsection{Canonical Symbol Key}
\begingroup
\small
\raggedright
\textbf{Source equation symbols.}\par
\begin{itemize}[leftmargin=1.35em,itemsep=1pt,topsep=1pt,parsep=0pt]
    \item $\mathcal{O}_{\mathrm{id0705}}$: canonical operator or carrier map for the entry.
    \item $\mathcal{X}_{\mathrm{id0705}}$: source domain or input state space for the entry map.
    \item $\mathcal{Y}_{\mathrm{id0705}}$: target domain or output state space for the entry map.
    \item $\mathcal{N}_{\mathrm{TRAWIN}}$: TRAWIN normalization operator applied to the entry relation.
    \item $\mathcal{Y}_{\mathrm{adm}}$: admissible output space selected by the source equation.
\end{itemize}
\endgroup
\end{minipage}\hfill
\begin{minipage}[t]{0.47\linewidth}
\tzpsection{Wolfram Form}
\begingroup
\scriptsize
\raggedright
\textbf{Source Wolfram model.}\par
\begin{Verbatim}[fontsize=\tiny,breaklines=true,breakanywhere=true]
(* TZPID ID0705 generated source Wolfram model *)
ClearAll[K0705, Cr0705, T, R, A, W, I, N];
eq0705$1 = HoldForm[Oid0705[Xid0705] -> Yid0705];
eq0705$2 = HoldForm[NTRAWIN[OId0705] == OId0705];
eq0705$3 = HoldForm[Element[OId0705[x], YAdm]];

K0705 = HoldForm[{eq0705$1, eq0705$2, eq0705$3}];

N = "normalized TRAWIN closure object for BERRY CONNECTION GAUGE INVARIANCE";
I = "BERRY CONNECTION GAUGE INVARIANCE integration/comparison layer";
W = "BERRY CONNECTION GAUGE INVARIANCE wave or operator carrier";
A = "BERRY CONNECTION GAUGE INVARIANCE admissible action/amplitude condition";
R = "BERRY CONNECTION GAUGE INVARIANCE rotation/curl preservation layer";
T = "BERRY CONNECTION GAUGE INVARIANCE local source condition";

Cr0705[expr_] := HoldForm[N[I[W[A[R[T[expr]]]]]]];

Result0705 = Cr0705[K0705];
\end{Verbatim}
\endgroup
\end{minipage}

\par\vspace{0.45\baselineskip}
\noindent\begin{minipage}[t]{\linewidth}
\begingroup
\small
\raggedright
\textbf{TRAWIN Operator Key.}\par
\noindent\textbf{Time/localization operator:} $\mathsf{T}\!\left[\mathrm{local}\right]$ BERRY CONNECTION GAUGE INVARIANCE local source condition.\par
\noindent\textbf{Rotation/curl operator:} $\mathsf{R}\!\left[\nabla\times\right]$ BERRY CONNECTION GAUGE INVARIANCE rotation/curl preservation layer.\par
\noindent\textbf{Action/amplitude operator:} $\mathsf{A}\!\left[\mathrm{admissible}\right]$ BERRY CONNECTION GAUGE INVARIANCE admissible action/amplitude condition.\par
\noindent\textbf{Wave/d'Alembertian operator:} $\mathsf{W}\!\left[\Box\right]$ BERRY CONNECTION GAUGE INVARIANCE wave or operator carrier.\par
\noindent\textbf{Integration operator:} $\mathsf{I}\!\left[\int\right]$ BERRY CONNECTION GAUGE INVARIANCE integration/comparison layer.\par
\noindent\textbf{Normalization operator:} $\mathsf{N}\!\left[\mathrm{normalized}\right]$ normalized TRAWIN closure object for BERRY CONNECTION GAUGE INVARIANCE.\par
\endgroup
\end{minipage}

\clearpage
\noindent\textbf{[ID 0705] BERRY CONNECTION GAUGE INVARIANCE}\par
\vspace{0.35em}
\tzpsection{Derivation Layer}
\paragraph{Primary Equation.}
\[
\mathcal{O}_{\mathrm{id0705}}:\mathcal{X}_{\mathrm{id0705}}\longrightarrow\mathcal{Y}_{\mathrm{id0705}}
\]
\[
\mathcal{N}_{\mathrm{TRAWIN}}(\mathcal{O}_{\mathrm{id0705}})=\mathcal{O}_{\mathrm{id0705}}
\]
\[
\mathcal{O}_{\mathrm{id0705}}(x)\in\mathcal{Y}_{\mathrm{adm}}
\]

\paragraph{Primary Derivation.}
The first line specifies a map, while the subsequent line provides its explicit form \( \mathcal{N}_{\mathrm{TRAWIN}}(\mathcal{O}_{\mathrm{id0705}})=\mathcal{O}_{\mathrm{id0705}} \). Verifying that the formula satisfies the declared mapping properties confirms that the map is well defined.
\paragraph{Equation Type.}
This statement is classified as a map relation.

\paragraph{Interpretation.}
The statement concerns the variables introduced in the displayed formula. Within the TRAWIN framework, the equation functions as the mapping carrier for the entry’s information content. Within the CHL view, the derivation provides an explicit term connecting the canonical proposition to its proof certificate.

\par\vspace{0.35em}
\noindent\textbf{Derivable Consequences.}\quad
\begingroup
\footnotesize
\raggedright
The canonical equation anchors BERRY CONNECTION GAUGE INVARIANCE by connecting the source condition to the derivation layer and proof certificate.
\par\endgroup

\tzpsection{Equation Support}
\noindent\textbf{Canonical Support Relation}\par
\[
\mathcal{O}_{\mathrm{id0705}}:\mathcal{X}_{\mathrm{id0705}}\longrightarrow\mathcal{Y}_{\mathrm{id0705}}
\]
\[
\mathcal{N}_{\mathrm{TRAWIN}}(\mathcal{O}_{\mathrm{id0705}})=\mathcal{O}_{\mathrm{id0705}}
\]
\[
\mathcal{O}_{\mathrm{id0705}}(x)\in\mathcal{Y}_{\mathrm{adm}}
\]
\noindent The support equation repeats the canonical relation so the derivation page remains self-contained.\par

\clearpage
\noindent\textbf{[ID 0705] BERRY CONNECTION GAUGE INVARIANCE}\par
\vspace{0.35em}
\tzpsection{Dictionary}
\begingroup\small
\tzpsection{Definition}
BERRY CONNECTION GAUGE INVARIANCE is a registry definition in Vortex and Cymatic Transport with secondary pressure from Resonance and Phase Locking.

% BEGIN TZP CONTEXTUAL MEANING
\tzpsection{Contextual Meaning}
\begingroup\footnotesize\setlength{\emergencystretch}{3em}\sloppy
\textbf{Formal statement:} A sub mu = ilangle psi | partial sub mu psi rangle implies oint mathbf A * dmathbf l = gamma. \textbf{Contextual meaning:} The entry reads BERRY CONNECTION GAUGE INVARIANCE as a controlled technical headword whose equation layer, proof route, and registry role are interpreted together. \textbf{Derivable consequences:} The entry now carries the actual connection-phase-gauge progression instead of malformed boldface export residue. \textbf{Lemma support:} Berry Connection; Geometric Phase; Gauge Shift.\par
\endgroup
% END TZP CONTEXTUAL MEANING

\tzpsection{Technical Interpretation}
Technically, BERRY CONNECTION GAUGE INVARIANCE is read through the local equation chain and the TRAWIN normalization recorded on the entry page.

\tzpsection{Equation Grounding}
Equation anchors: A sub mu =i psi mid partial sub mu psi; gamma=oint A dot d l; and A sub mu mapsto A sub mu +partial sub mu chi,; gammamapsto gamma bmod 2pi. Proof route: show that the Berry connection integrates to the stated geometric phase and remains physically invariant under the stated gauge shift. Formalization route: prove that the connection one-form yields the closed-cycle phase modulo the stated gauge transformation.

\tzpsection{Interpretive Role}
The entry now carries the actual connection-phase-gauge progression instead of malformed boldface export residue.
\endgroup

\tzpsection{TZP Thesaurus}
\begin{enumerate}[leftmargin=1.6em,itemsep=6pt,topsep=4pt]
\item \textbf{Arbitrary Phase-Shift Immunity}: The beautiful mathematical proof that even if you completely change the completely made-up, invisible labels you use to track the quantum wave, the actual, physical reality of the knot stays exactly the same.
\item \textbf{Coordinate-Free Topological Truth}: The realization that the deep, knotted magic of the system is a fundamental, undeniable reality of the universe, and not just a weird side effect of how you chose to write down the equations.
\item \textbf{Subjective-Observation Independence}: The comforting guarantee that the incredible, bizarre twisting of the geometric phase doesn't care at all about who is looking at it, or what math they are using to measure it.
\end{enumerate}

\tzpsection{Encyclopedia Mathematica}
\begingroup\footnotesize\sloppy
The visual plate below is reserved for the source-specific equation and progression graphic for [ID 0705]. It is included only as a companion to the canonical equation, not as a substitute for the formal text.
\par\endgroup
\begingroup\footnotesize\itshape
Visual plate pending source-specific approval for [ID 0705].
\par\endgroup

\clearpage
\begingroup\tzpsupportsetup
\noindent\textbf{\fontsize{10.8pt}{12.8pt}\selectfont [ID 0705] Constructive Logic and Proof Certificate}\par
\noindent\textbf{\fontsize{10pt}{12pt}\selectfont BERRY CONNECTION GAUGE INVARIANCE}\par
\vspace{0.45em}
\begingroup
\footnotesize
\setlength{\parskip}{0.22em}
\setlength{\parindent}{0pt}
\noindent\textbf{Curry--Howard--Lambek Interpretation}\par
Under the Curry--Howard--Lambek correspondence, this registry entry is read as a typed proof
object: the stated source condition supplies the proposition, the derivation supplies the
morphism, and the proof certificate records the machine handles needed to check the obligation
in the formal registry.

\vspace{0.45em}
\noindent\textbf{CHL Proof}\par
\textbf{Statement:} BERRY CONNECTION GAUGE INVARIANCE is read as a source-backed proof obligation.

\textbf{Logic Validation:}\\
\textbf{Proposition:} ``BERRY CONNECTION GAUGE INVARIANCE is admissible under the named source equation and derivation.''\\
\textbf{Proof:} The canonical statement supplies the proposition, the derivation supplies the witness, and the TRAWIN operator chain records the normalization path.

\textbf{VERDICT:} \ensuremath{\checkmark}\ \textbf{TRUE} --- conditional registry validation

\textbf{Formal Status:} \ensuremath{\checkmark}\ THEORETICAL -- source-backed CHL proof obligation

\textbf{Physical Status:} THEORETICAL -- requires model-specific empirical interpretation
\par\vspace{0.45em}
\noindent{\tiny\textbf{Isabelle/HOL.} \tzpplaincode{registry\_number\_matches, equation\_count\_matches}}\par
\vfill
\begingroup\footnotesize\itshape
Proof visual pending source-specific approval for [ID 0705].
\par\endgroup
\vfill
\endgroup
\endgroup
